% Vietnamese translation for OI Public Library
% Source-PDF: IOI中国国家候选队论文/2002/周文超.pdf
\documentclass[11pt]{article}
\usepackage{oipl-vn}

\title{Ứng dụng cấu trúc cây trong lập trình}
\author{Zhou Wenchao}
\date{}

\begin{document}
\maketitle

\origtitle{Applications of tree structures in programming}
\sourcepath{IOI China candidate team papers / 2002 / Zhou Wenchao.pdf}

\renewcommand{\contentsname}{Mục lục}
\tableofcontents

\section{Dẫn nhập}

Trong những năm gần đây, vì nhiều kỳ thi lập trình đã chuyển sang dùng
Free Pascal, yêu cầu về hiệu quả bộ nhớ đối với thuật toán được nới lỏng hơn,
trong khi yêu cầu về thời gian chạy lại cao hơn. Điều đó đòi hỏi thí sinh
không chỉ nắm chắc các thuật toán thông thường, mà còn phải mạnh dạn cải tiến
và xây dựng thuật toán hiệu quả hơn để giải quyết bài toán.

Trong lập trình trước đây, cấu trúc liên kết được dùng khá nhiều. Cấu trúc
liên kết thật sự có ưu điểm là dễ lập trình và dễ hiểu, nhưng nó có một điểm
yếu chí mạng: quan hệ giữa hai phần tử kề nhau không được thể hiện rõ ràng.
Cấu trúc cây lại làm được điều này rất tốt. Bài viết giới thiệu ba ví dụ tiêu
biểu: tập hợp rời nhau, cây đoạn và mảng cây.

\section{Tập hợp rời nhau}

Trong thi đấu, ta thường gặp dạng bài: cho các quan hệ giữa các phần tử, hãy
chia các phần tử thành một số tập sao cho các phần tử trong cùng một tập có
liên hệ trực tiếp hoặc gián tiếp với nhau. Dạng bài này chủ yếu cần hai thao
tác: hợp nhất hai tập và tìm tập chứa một phần tử. Vì vậy cấu trúc dữ liệu
tương ứng thường được gọi là \term{tập hợp rời nhau}, hay \term{union--find}.

\subsection{Cách lưu bằng danh sách liên kết}

Danh sách liên kết từng được dùng phổ biến để hiện thực union--find. Mỗi phần
tử trong danh sách có hai con trỏ: một con trỏ trỏ tới phần tử kế tiếp trong
cùng tập, con trỏ còn lại trỏ tới phần tử đầu danh sách. Với cách lưu này,
thao tác tìm tập có độ phức tạp chỉ \(O(1)\), vì từ một phần tử có thể đến
ngay đại diện của tập. Tuy nhiên, thao tác hợp nhất hai tập có thể phải đổi
con trỏ đầu danh sách của mọi phần tử trong một tập, nên độ phức tạp lên đến
\(O(n)\). Nếu muốn cả hai thao tác cơ bản đều nhanh, cấu trúc liên kết không
còn phù hợp.

\subsection{Cách lưu bằng cây}

Cấu trúc cây có thể hỗ trợ union--find tốt hơn. Khác với cây thông thường,
mỗi đỉnh không ghi danh sách các con mà ghi cha của nó. Gốc của cây chính là
đại diện của tập. Với mảng \(\texttt{parent}\), ta có:
\[
  \texttt{parent}[x]=x \quad\Longleftrightarrow\quad x
  \text{ là gốc của cây.}
\]

Hai thao tác cơ bản có thể viết như sau.

\begin{verbatim}
function FIND(x)
begin
    while parent[x] <> x do x <- parent[x]
    return x
end

procedure UNION(x, y)
begin
    rx <- FIND(x)
    ry <- FIND(y)
    if rx <> ry then parent[rx] <- ry
end
\end{verbatim}

Việc hợp nhất hai tập rất đơn giản: chỉ cần nối gốc của một cây vào gốc của
cây kia, thao tác nối này mất \(O(1)\) sau khi đã tìm được hai gốc. Do đó
thời gian của union--find phụ thuộc chủ yếu vào tốc độ tìm gốc. Nếu truy vấn
trực tiếp trong cây, thời gian tìm tuyến tính theo độ sâu của cây. Trung bình
có thể xem là \(O(\log N)\) khi cây khá cân bằng, nhưng trong trường hợp xấu
nhất cây suy biến thành một dây, làm mỗi lần tìm tốn \(O(N)\).

\subsection{Nén đường đi}

Ta còn có thể giảm độ phức tạp của thao tác tìm bằng \term{nén đường đi}. Ý
tưởng rất đơn giản: trong lúc tìm gốc của một phần tử, đồng thời làm giảm độ
sâu của các đỉnh nằm trên đường đi đó bằng cách cho chúng trỏ thẳng tới gốc.

\begin{verbatim}
function FIND(x)
begin
    if parent[x] <> x then parent[x] <- FIND(parent[x])
    return parent[x]
end
\end{verbatim}

Sau khi dùng nén đường đi, mỗi truy vấn có độ phức tạp khấu hao là hàm ngược
của hàm Ackermann, ký hiệu \(\alpha(n)\), một hàm tăng cực kỳ chậm. Với mọi
giá trị \(n\) có thể gặp trong thực tế, \(\alpha(n)\) không vượt quá \(5\).

So với cách lưu bằng danh sách liên kết, cách lưu bằng cây có ưu thế rất rõ:
cài đặt vẫn ngắn gọn, trong khi hiệu quả thời gian cao hơn nhiều.

\section{Cây đoạn}

Khi xử lý các bài toán liên quan đến diện tích, chu vi và các đại lượng hình
học tương tự, ta thường không cần kiến thức toán học quá sâu; khó khăn nằm ở
việc nâng cao hiệu quả xử lý. Muốn vậy, nhiều khi phải thay đổi nền tảng của
thuật toán, tức thay đổi cấu trúc dữ liệu. Cây đoạn là một cấu trúc dữ liệu
đặc biệt phục vụ mục tiêu đó.

Ta xét bài toán cơ sở: cho \(n\) đoạn thẳng trên một trục, hãy xác định độ dài
phần được các đoạn này phủ. Qua bài toán này, ta lần lượt thấy định nghĩa cây
đoạn, thao tác thêm/xóa đoạn, và cách tính độ đo cùng số đoạn liên tục.

\subsection{Định nghĩa}

Cây đoạn là một cây nhị phân chia một khoảng thành các khoảng đơn vị
\([i,i+1]\). Mỗi khoảng đơn vị tương ứng với một lá trong cây. Mỗi nút biểu
diễn một khoảng \([l,r]\), và có một biến \(\texttt{count}\) ghi số đoạn đang
phủ hoàn toàn khoảng của nút đó.

Ví dụ, cây đoạn của khoảng \([1,7]\) có thể chia như sau. Trong ví dụ có một
đoạn \([3,6]\); các khoảng sơ cấp bị đoạn này phủ là \([3,4]\), \([4,5]\) và
\([5,6]\).

\begin{verbatim}
                         [1,7]
                    /             \
                [1,4]             [4,7]
              /      \           /     \
          [1,2]      [2,4]   [4,5]   [5,7]
                    /   \            /   \
                [2,3] [3,4]      [5,6] [6,7]
\end{verbatim}

Đoạn \([3,6]\) phủ các lá \([3,4]\), \([4,5]\) và \([5,6]\).

\subsection{Thêm và xóa đoạn}

Thao tác thêm và xóa đoạn trong cây đoạn giống nhau về cấu trúc. Ta quét đệ
quy xuống các con cho đến khi đoạn cần xử lý phủ hoàn toàn khoảng của nút
hiện tại; khi đó chỉ cập nhật \(\texttt{count}\) tại nút này.

\begin{verbatim}
procedure ADD(v, L, R, delta)
begin
    if L <= left[v] and right[v] <= R then
        count[v] <- count[v] + delta
    else
    begin
        mid <- (left[v] + right[v]) div 2
        if L < mid then ADD(leftChild[v], L, R, delta)
        if mid < R then ADD(rightChild[v], L, R, delta)
        PULL(v)
    end
end
\end{verbatim}

Khi \(\texttt{delta}=1\), thủ tục là thêm đoạn; khi \(\texttt{delta}=-1\),
thủ tục là xóa đoạn. Phân tích cho thấy thời gian thêm hoặc xóa một đoạn đều
là \(O(\log N)\), với \(N\) là số khoảng đơn vị hoặc số điểm chia sau khi rời
rạc hóa.

\subsection{Độ đo}

Độ đo \(m\) của một nút là tổng độ dài trong khoảng của nút đang được phủ bởi
ít nhất một đoạn. Nếu nút biểu diễn khoảng \([i,j]\), ta có:
\[
m =
\begin{cases}
  j-i, & \text{nếu } \texttt{count}>0,\\
  0, & \text{nếu } \texttt{count}=0 \text{ và nút là lá},\\
  m_{\mathrm{left}}+m_{\mathrm{right}},
    & \text{nếu } \texttt{count}=0 \text{ và nút là nút trong}.
\end{cases}
\]

Nói cách khác, nếu chính khoảng của nút đã bị một đoạn phủ hoàn toàn, độ đo
của nút là toàn bộ chiều dài khoảng. Nếu không, độ đo phải được lấy từ hai
cây con.

\subsection{Số đoạn liên tục}

Số đoạn liên tục, ký hiệu \(\texttt{line}\), là số thành phần liên thông trong
hợp các đoạn nằm trong khoảng của nút. Đại lượng này không thể tính bằng cách
cộng đơn giản số đoạn liên tục của hai con, vì một đoạn ở con trái và một
đoạn ở con phải có thể nối liền nhau tại ranh giới.

Do đó ta thêm hai biến. Biến biên trái là
\[
\texttt{lbd}=
\begin{cases}
1, & \text{nếu đầu trái của khoảng được phủ},\\
0, & \text{nếu đầu trái của khoảng không được phủ}.
\end{cases}
\]
Biến biên phải là
\[
\texttt{rbd}=
\begin{cases}
1, & \text{nếu đầu phải của khoảng được phủ},\\
0, & \text{nếu đầu phải của khoảng không được phủ}.
\end{cases}
\]

Khi nút biểu diễn khoảng \([i,j]\), ta cập nhật \(\texttt{line}\) như sau:
\[
\texttt{line} =
\begin{cases}
  1, & \text{nếu } \texttt{count}>0,\\
  0, & \text{nếu } \texttt{count}=0 \text{ và nút là lá},\\
  \texttt{line}_{L}+\texttt{line}_{R}-1,
    & \begin{array}{l}
      \text{nếu } \texttt{count}=0,\ \text{nút là nút trong,}\\
      \texttt{rbd}_{L}=1 \text{ và } \texttt{lbd}_{R}=1,
    \end{array}\\
  \texttt{line}_{L}+\texttt{line}_{R},
    & \text{trong trường hợp còn lại.}
\end{cases}
\]

Trường hợp trừ \(1\) là trường hợp hai thành phần ở hai cây con chạm nhau tại
điểm chia, vì đầu phải của con trái và đầu trái của con phải đều đang được
phủ. Các biến biên cũng được cập nhật theo cùng nguyên tắc: nếu
\(\texttt{count}>0\), cả hai biên đều bằng \(1\); nếu không thì lấy biên
ngoài tương ứng từ hai con.

\section{Mảng cây}

Bài \emph{Mobile} của IOI 2001 đã gây khó khăn cho nhiều thí sinh. Nội dung
bài toán rất đơn giản: trong một ma trận, ta thay đổi giá trị phần tử để cập
nhật trạng thái ma trận, đồng thời cần tính tổng các phần tử trong một ma
trận con. Điểm khó nằm ở quy mô dữ liệu rất lớn. Phần này giới thiệu một cấu
trúc dữ liệu mới: \term{mảng cây}, tức \term{Fenwick tree}.

Nhờ mảng cây, độ phức tạp lập trình tăng lên một chút, nhưng hiệu quả thời
gian của chương trình tăng đáng kể. Tư tưởng chính vẫn là lợi dụng cấu trúc
cây để thu hẹp phạm vi tìm kiếm và gom thông tin, làm cho mỗi lần cập nhật
hoặc tính tổng chỉ liên quan đến ít biến nhất có thể.

\subsection{Từ tổng tiền tố một chiều}

Trước hết đơn giản hóa bài toán về một chiều. Cho dãy
\[
  a[1],a[2],\ldots,a[n].
\]

Cách thứ nhất là tính trực tiếp trên dãy gốc. Cập nhật một phần tử mất
\(O(1)\), nhưng trong trường hợp xấu nhất, tính tổng một dãy con mất \(O(n)\).

Cách thứ hai là thêm mảng tiền tố \(b\), trong đó
\[
  b[i]=a[1]+a[2]+\cdots+a[i].
\]
Khi đó truy vấn tổng dãy con có thể thực hiện trong \(O(1)\), nhưng thay đổi
\(a[i]\) sẽ ảnh hưởng đến \(b[i],b[i+1],\ldots,b[n]\), nên cập nhật có thể
tốn \(O(n)\).

Hai cách trên đều làm một trong hai thao tác bị chậm. Ta cần một cách cân
bằng hơn.

\subsection{Xây dựng mảng C}

Ta thêm mảng \(C\), trong đó
\[
  C[i]=a[i-2^k+1]+a[i-2^k+2]+\cdots+a[i],
\]
với \(k\) là số chữ số \(0\) liên tiếp ở cuối biểu diễn nhị phân của \(i\).
Tương đương, \(2^k=\texttt{lowbit}(i)\), nên \(C[i]\) lưu tổng
\(\texttt{lowbit}(i)\) phần tử kết thúc tại vị trí \(i\).

Từ định nghĩa của mảng \(C\), ta có:
\[
\begin{aligned}
C[1]&=a[1],\\
C[2]&=a[1]+a[2]=C[1]+a[2],\\
C[3]&=a[3],\\
C[4]&=a[1]+a[2]+a[3]+a[4]=C[2]+C[3]+a[4],\\
C[5]&=a[5],\\
C[6]&=a[5]+a[6]=C[5]+a[6].
\end{aligned}
\]

Cấu trúc bao phủ của mảng \(C\) tương ứng với một cây, nên được gọi là mảng
cây. Sau khi phân tích hai thao tác cập nhật và tính tổng dãy con, ta thấy
cả hai đều có độ phức tạp \(O(\log N)\), cải thiện rõ rệt so với hai cách cơ
sở.

\subsection{Cập nhật giá trị phần tử}

\paragraph{Định lý.}
Giả sử khi thay đổi \(a[k]\), các phần tử của mảng \(C\) bị ảnh hưởng là
\[
  C[p_1], C[p_2], \ldots, C[p_m].
\]
Khi đó
\[
  p_1=k,\qquad p_{i+1}=p_i+2^{l_i},
\]
trong đó \(l_i\) là số chữ số \(0\) liên tiếp ở cuối biểu diễn nhị phân của
\(p_i\). Nói cách khác, \(2^{l_i}=\texttt{lowbit}(p_i)\).

Từ định lý này, ta có cách cập nhật phần tử: nếu cộng \(x\) vào \(a[k]\), thì
các vị trí \(C[p_1],C[p_2],\ldots,C[p_m]\), với \(p_m\le n<p_{m+1}\), cũng
phải được cộng \(x\).

\begin{verbatim}
procedure ADD(k, x)
begin
    while k <= n do
    begin
        C[k] <- C[k] + x
        k <- k + LOWBIT(k)
    end
end
\end{verbatim}

Ví dụ, với dãy \(a[1],\ldots,a[9]\), nếu cộng \(x\) vào \(a[3]\), ta có
\[
  p_1=3,\qquad
  p_2=3+2^0=4,\qquad
  p_3=4+2^2=8,\qquad
  p_4=8+2^3=16>9.
\]
Vì vậy cần cộng \(x\) vào \(C[3]\), \(C[4]\) và \(C[8]\).

\paragraph{Bổ đề.}
Nếu dãy vị trí bị \(a[k]\) ảnh hưởng là
\[
  C[p_1],C[p_2],\ldots,C[p_m],
  \qquad p_1<p_2<\cdots<p_m,
\]
và \(l_i\) là số chữ số \(0\) liên tiếp ở cuối biểu diễn nhị phân của
\(p_i\), thì
\[
  l_1<l_2<\cdots<l_m.
\]

\paragraph{Chứng minh.}
Vì \(C[p_i]\) chứa \(a[k]\), nên
\[
  p_i-2^{l_i}+1 \le k \le p_i.
\]
Giả sử tồn tại \(i\) sao cho \(l_i\ge l_{i+1}\). Từ việc \(C[p_{i+1}]\) cũng
chứa \(a[k]\), ta có
\[
  p_{i+1}-2^{l_{i+1}}+1 \le k \le p_i,
\]
suy ra
\[
  p_{i+1}\le p_i+2^{l_{i+1}}-1. \tag{1}
\]
Mặt khác, do \(p_i<p_{i+1}\) và \(l_i\ge l_{i+1}\), vị trí tiếp theo có thể
có ít nhất \(l_{i+1}\) chữ số \(0\) ở cuối không thể xuất hiện trước
\[
  p_i+2^{l_{i+1}}.
\]
Vì thế
\[
  p_{i+1}\ge p_i+2^{l_{i+1}}. \tag{2}
\]
Hai bất đẳng thức (1) và (2) mâu thuẫn. Do đó
\(l_1<l_2<\cdots<l_m\).

\paragraph{Chứng minh định lý.}
Vì \(p_1<p_2<\cdots<p_m\) và mọi \(C[p_i]\) đều chứa \(a[k]\), vị trí nhỏ
nhất chắc chắn là \(p_1=k\). Trong dãy \(p\), vị trí
\[
  p_i+2^{l_i}
\]
là vị trí nhỏ nhất sau \(p_i\) có số chữ số \(0\) ở cuối lớn hơn \(l_i\). Nếu
có một \(p_i+x\) nhỏ hơn vị trí này, thì \(x<2^{l_i}\), số chữ số nhị phân
cần xét của \(x\) không đủ để tạo ra nhiều chữ số \(0\) ở cuối hơn, mâu thuẫn
với bổ đề.

Đặt \(p_{i+1}=p_i+2^{l_i}\). Khi đó \(p_{i+1}\) có ít nhất \(l_i+1\) chữ số
\(0\) ở cuối. Từ
\[
  p_i-2^{l_i}+1\le k\le p_i
\]
suy ra
\[
  p_{i+1}-2^{l_{i+1}}+1
  \le p_i+2^{l_i}-2\cdot 2^{l_i}+1
  = p_i-2^{l_i}+1
  \le k \le p_i \le p_{i+1}.
\]
Vì thế \(C[p_{i+1}]\) cũng chứa \(a[k]\), và quan hệ truy hồi cần chứng minh là
\[
  p_{i+1}=p_i+2^{l_i}.
\]

\subsection{Tính tổng dãy con}

Tổng một dãy con có thể chuyển thành hiệu của hai tổng tiền tố, nên chỉ cần
biết cách tính
\[
  S(k)=a[1]+a[2]+\cdots+a[k].
\]

Theo định nghĩa
\[
  C[k]=a[k-2^l+1]+\cdots+a[k],
\]
với \(l\) là số chữ số \(0\) liên tiếp ở cuối biểu diễn nhị phân của \(k\).
Bắt đầu từ \(k_1=k\), ta dùng công thức
\[
  k_{i+1}=k_i-2^{l_i},
\]
trong đó \(l_i\) là số chữ số \(0\) liên tiếp ở cuối biểu diễn nhị phân của
\(k_i\). Tiếp tục cho đến khi \(k_{m+1}=0\), ta được
\[
  S(k)=C[k_1]+C[k_2]+\cdots+C[k_m].
\]

\begin{verbatim}
function SUM(k)
begin
    ans <- 0
    while k > 0 do
    begin
        ans <- ans + C[k]
        k <- k - LOWBIT(k)
    end
    return ans
end
\end{verbatim}

Ví dụ, để tính
\[
  a[1]+a[2]+a[3]+a[4]+a[5]+a[6]+a[7],
\]
ta có
\[
  k_1=7,\qquad
  k_2=7-2^0=6,\qquad
  k_3=6-2^1=4,\qquad
  k_4=4-2^2=0.
\]
Do đó
\[
  a[1]+a[2]+a[3]+a[4]+a[5]+a[6]+a[7]
  =C[7]+C[6]+C[4].
\]

\subsection{Tổng ma trận con}

Mảng cây cũng có thể mở rộng sang hai chiều. Khi đó ta chỉ cần tính tổng của
hình chữ nhật bắt đầu từ \((1,1)\) đến \((x,y)\), rồi dùng cộng trừ để lấy
tổng của một ma trận con bất kỳ.

Gọi \(S(x,y)\) là tổng hình chữ nhật \([1,x]\times[1,y]\). Khi cần tổng của
\([x_1,x_2]\times[y_1,y_2]\), ta dùng:
\[
\begin{aligned}
\operatorname{sum}(x_1,y_1,x_2,y_2)
  ={}& S(x_2,y_2)-S(x_1-1,y_2)\\
     &-S(x_2,y_1-1)+S(x_1-1,y_1-1).
\end{aligned}
\]

Cập nhật và truy vấn hai chiều có dạng lồng hai vòng lặp của phiên bản một
chiều:

\begin{verbatim}
procedure ADD2D(x, y, delta)
begin
    i <- x
    while i <= n do
    begin
        j <- y
        while j <= m do
        begin
            C[i][j] <- C[i][j] + delta
            j <- j + LOWBIT(j)
        end
        i <- i + LOWBIT(i)
    end
end

function SUM2D(x, y)
begin
    ans <- 0
    i <- x
    while i > 0 do
    begin
        j <- y
        while j > 0 do
        begin
            ans <- ans + C[i][j]
            j <- j - LOWBIT(j)
        end
        i <- i - LOWBIT(i)
    end
    return ans
end
\end{verbatim}

Mỗi thao tác trong hai chiều có độ phức tạp \(O(\log n\log m)\).

\section{Kết luận}

Qua các ví dụ trên, ta thấy khi dùng cấu trúc cây để xử lý bài toán, phạm vi
tìm kiếm thường được thu hẹp đáng kể. Trong mỗi lần truy vấn, thuật toán có
thể chọn đúng một đường đi bắt đầu từ gốc, tránh được các phép tìm kiếm không
cần thiết.

Vì thông thường ta chỉ đi theo một đường từ gốc xuống, tốc độ duyệt cây phụ
thuộc rất nhiều vào độ sâu của cây. Do đó, làm thế nào để giảm chiều cao của
cây trở thành điểm then chốt của thuật toán.

So với cây theo nghĩa truyền thống, mảng cây nhìn bề ngoài không hẳn là một
cây. Tuy vậy, tư tưởng cốt lõi của nó vẫn cùng một bản chất với cấu trúc cây:
tổ chức thông tin theo từng mức để thu hẹp phạm vi cập nhật và truy vấn.

\end{document}
